\documentclass[a4paper,12pt]{article}
\usepackage[utf8]{vietnam}
\usepackage[left=2cm,right=2cm,top=2cm,bottom=2cm]{geometry}
\usepackage{amsmath,amssymb}
\usepackage{ex_test}

\begin{document}

\begin{center}
%Tiêu đề
% Cách dùng: \tieude{tổng số trang}{mã đề}{tiêu đề riêng (để trống nếu dùng mặc định)}{role: hs hoặc gv}{lop: 10/11/12}
\newcommand{\tieude}[5]{
	\def\customtitle{#3}
	\ifthenelse{\equal{\customtitle}{}}
		{\def\thetitle{ÔN TẬP KIỂM TRA}}
		{\def\thetitle{\customtitle}}
	\noindent
	%Trái
	\begin{minipage}[b]{8.5cm}
		\centerline{\textbf{\fontsize{11}{0}\selectfont \hspace{2cm} TRƯỜNG THPT CHUYÊN HÙNG VƯƠNG}}
		\centerline{(\textit{Đề thi có #1\ trang})}
	\end{minipage}\hspace{1.5cm}
	%Phải
	\begin{minipage}[b]{9cm}
		\centerline{\textbf{\fontsize{11}{0}\selectfont \thetitle}}
		\centerline{\textbf{\fontsize{11}{0}\selectfont MÔN TOÁN, LỚP #5 (KHÔNG CHUYÊN)}}
	\end{minipage}\\
	\vspace{0.4cm}
	\ifthenelse{\equal{#4}{gv}}
		{\rightline{
			\setlength\fboxrule{1pt}
			\setlength\fboxsep{4pt}
			\fbox{\bf Mã đề #2}
		}}
		{}
	\vspace{4pt}
}
\end{center}

\begin{ex}%%[?]
Với $n \in \mathbb{N}$.
\choiceTFn[4]{Mệnh đề ``$n$ là một số không âm'' là một mệnh đề sai}
{\True Mệnh đề ``$\exists n \in \mathbb{N},\, 2 \nmid (n-1)n$'' là mệnh đề sai}
{\True Mệnh đề ``$\exists n \in \mathbb{N}$ sao cho $(n-1)n$ là hợp số'' là mệnh đề đúng}
{\True Mệnh đề ``$\exists n \in \mathbb{N},\, 9 \nmid (n^3-n)$'' là mệnh đề đúng}

\loigiai{
\begin{itemchoice}
\itemch Sai. $\mathbb{N} = \{0;1;2;\ldots\}$ nên mọi $n \in \mathbb{N}$ đều thỏa $n \geq 0$.
\itemch Đúng. Tích hai số nguyên liên tiếp luôn là số chẵn.
\itemch Đúng. $n=3$: $2\cdot3=6$ là hợp số.
\itemch Đúng. Với $n=2$ thì $n^3-n=6$, mà $6$ không chia hết cho $9$.
\end{itemchoice}
}
\end{ex}

\begin{center}\fbox{\textbf{[THIẾU CÂU HỎI: L10\_C1\_B2\_NB017 - trac\_nghiem]}}\end{center}

\begin{center}{\small\textit{Lỗi khi chạy hàm sinh câu (AttributeError): '\_cython\_3\_2\_8.cython\_function\_or\_method' object has no attribute 'randint'}}\end{center}

\begin{ex}%%[?]
Phát biểu bằng lời mệnh đề:
``$\forall x\in \mathbb{R}, x^2 \leq 0$''.
\choice{Mọi số thực đều có bình phương khác $0$}
{Tồn tại số thực mà bình phương của nó lớn hơn $0$}
{Mọi số thực đều có bình phương lớn hơn $0$}
{\True Mọi số thực đều có bình phương nhỏ hơn hoặc bằng $0$}


\loigiai{
Kí hiệu $\forall$ phát biểu là "mọi",
kí hiệu $\exists$ phát biểu là "tồn tại".
}
\end{ex}

\begin{ex}%%[?]
Cho hai mệnh đề sau:\\
            $P \colon$ ``Tập thể dục thường xuyên'';\\
            $Q \colon$ ``Sức khỏe được cải thiện''.\\
            Hãy phát biểu mệnh đề kéo theo $P \Rightarrow Q$.
\choice{\True Nếu tập thể dục thường xuyên thì sức khỏe được cải thiện}
{Nếu sức khỏe được cải thiện thì tập thể dục thường xuyên}
{Nếu tập thể dục thường xuyên thì không có chuyện sức khỏe được cải thiện}
{Tập thể dục thường xuyên khi và chỉ khi sức khỏe được cải thiện}


\loigiai{
Mệnh đề kéo theo $P \Rightarrow Q$ được phát biểu dưới dạng: ``Nếu $P$ thì $Q$''. Do đó đáp án đúng là: ``Nếu tập thể dục thường xuyên thì sức khỏe được cải thiện''.
}
\end{ex}

\begin{center}\fbox{\textbf{[THIẾU CÂU HỎI: L10\_C1\_B1\_NB010 - trac\_nghiem]}}\end{center}

\begin{center}{\small\textit{Lỗi khi chạy hàm sinh câu (AttributeError): '\_cython\_3\_2\_8.cython\_function\_or\_method' object has no attribute 'randint'}}\end{center}

\begin{ex}%%[?]
Trong các khẳng định sau, khẳng định nào đúng?
\choice{\True $\exists x\in\mathbb R,\ 9x+(-14)=0$}
{$\forall x\in\mathbb R,\ 9x+(-14)=0$}
{$\nexists x\in\mathbb R,\ 9x+(-14)=0$}
{$\forall x\in\mathbb R,\ 9x+(-14)>0$}


\loigiai{
Khẳng định đúng là phương án đã chọn.
}
\end{ex}

\begin{center}\fbox{\textbf{[THIẾU CÂU HỎI: L10\_C1\_B2\_VD020 - trac\_nghiem]}}\end{center}

\begin{center}{\small\textit{Lỗi khi chạy hàm sinh câu (AttributeError): '\_cython\_3\_2\_8.cython\_function\_or\_method' object has no attribute 'choice'}}\end{center}

\begin{center}\fbox{\textbf{[THIẾU CÂU HỎI: L10\_C1\_B2\_VD020 - tra\_loi\_ngan]}}\end{center}

\begin{center}{\small\textit{L10\_C1\_B2\_VD020 (SA): không có Generator nào khớp trong Mapping.}}\end{center}

\begin{center}\fbox{\textbf{[THIẾU CÂU HỎI: L10\_C1\_B2\_VD021 - tra\_loi\_ngan]}}\end{center}

\begin{center}{\small\textit{L10\_C1\_B2\_VD021 (SA): không có Generator nào khớp trong Mapping.}}\end{center}

\begin{center}\fbox{\textbf{[THIẾU CÂU HỎI: L10\_C1\_B2\_VD021 - tu\_luan]}}\end{center}

\begin{center}{\small\textit{L10\_C1\_B2\_VD021 (TL): không có Generator nào khớp trong Mapping.}}\end{center}

\begin{ex}%%[?]
Tại một sự kiện có 100 người tham gia khảo sát về hai sản phẩm A và B. Biết có 61 người chọn sản phẩm A, 60 người chọn sản phẩm B và 51 người chọn cả hai sản phẩm. Tính: 
\begin{listEX}[1]
\item Số người đã chọn ít nhất một sản phẩm? \SA[4]{$70$}
\item Số người không chọn sản phẩm nào? \SA[4]{$30$}
\end{listEX}

 \loigiai{
\begin{listEX}[1]
\item Theo nguyên lý bù trừ, số người chọn ít nhất một sản phẩm là: 61 + 60 - 51 = 70 (người).
\item Số người không chọn sản phẩm nào là: 100 - 70 = 30 (người).
\end{listEX}

}
\end{ex}

\begin{ex}%%[?]
Tại một sự kiện có 100 người tham gia khảo sát về hai sản phẩm A và B. Biết có 59 người chọn sản phẩm A, 52 người chọn sản phẩm B và 33 người chọn cả hai sản phẩm. Tính: 
\begin{listEX}[1]
\item Số người đã chọn ít nhất một sản phẩm? \SA[4]{$78$}
\item Số người không chọn sản phẩm nào? \SA[4]{$22$}
\end{listEX}

 \loigiai{
\begin{listEX}[1]
\item Theo nguyên lý bù trừ, số người chọn ít nhất một sản phẩm là: 59 + 52 - 33 = 78 (người).
\item Số người không chọn sản phẩm nào là: 100 - 78 = 22 (người).
\end{listEX}

}
\end{ex}



\end{document}